\section{Results: In-distribution}

\subsection{Training dynamics}

Both networks converged stably under the early-stopping regime, and the learning curves in \ref{fig:curves} show no evidence of overfitting: training and validation losses descend in step, and the validation accuracy plateaus rather than degrading after its peak. The teacher required 71 epochs in total, comprising 30 epochs with the backbone frozen followed by 41 epochs of whole-network fine-tuning before early stopping, and reached a final validation accuracy of $0.9442$. The distilled student converged considerably faster, early-stopping after 34 epochs at a final validation accuracy of $0.9661$. The absence of a widening train--validation gap in either curve indicates that the weighted cross-entropy objective for the teacher and the Hinton distillation objective for the student, combined with the augmentation pipeline, were sufficient to regularize models of $18.5$~M and $4.67$~M parameters, respectively, on the $4{,}012$-image training split.

\subsection{Classification performance}

On the held-out SIPaKMeD validation set ($n=1{,}003$), the $4.67$~M-parameter student does not merely approximate its $18.5$~M-parameter teacher; it surpasses it on every metric we report (\ref{tab:indist}). The student attains an accuracy of $0.9661$, balanced accuracy of $0.9680$, macro-F1 of $0.9671$, Cohen's $\kappa$ of $0.9576$, and a one-vs-rest macro AUROC of $0.9981$, against the teacher's $0.9432$, $0.9458$, $0.9449$, $0.9289$, and $0.9955$. That a model with roughly one-quarter of the teacher's parameters improves on it across accuracy, agreement, and ranking metrics is consistent with the regularizing effect of knowledge distillation, whereby the soft teacher targets convey inter-class similarity structure that a hard one-hot objective discards, answering RQ1 affirmatively: the distilled student matches and indeed exceeds the teacher at a fraction of the parameter budget.

The per-class breakdown (\ref{tab:perclass}) and the confusion matrix in \ref{fig:cm_student} locate the student's residual error. Performance is near-saturated on the two morphologically distinctive normal classes, with F1 scores of $0.995$ for Superficial-Intermediate and $0.994$ for Parabasal, and remains strong for Dyskeratotic ($0.971$) and Metaplastic ($0.957$). The clear weak point is the Koilocytotic class, whose F1 of $0.919$ is driven by a recall of only $0.873$: roughly one in eight koilocytotic cells is misassigned, predominantly to its morphological neighbours. This is cytologically unsurprising, as koilocytotic change---perinuclear haloing and nuclear atypia of intermediate severity---occupies a continuum between the metaplastic and dyskeratotic appearances rather than a sharply separable region of feature space, so the confusions concentrate exactly where human graders also disagree. Notably, Koilocytotic is also the teacher's weakest class (F1 $0.877$), and the student's improvement here ($+0.042$ F1) accounts for much of its overall margin over the teacher.

\subsection{Calibration}

Strong discrimination is matched by faithful confidence estimates. The student's $15$-bin Expected Calibration Error is $0.0152$, and the reliability diagram in \ref{fig:reliability} tracks the identity line closely across the confidence range, indicating that predicted probabilities can be read as approximately trustworthy likelihoods of correctness. The teacher is marginally better calibrated still, with an ECE of $0.0131$, so the small calibration cost of distillation is more than offset by the student's gains in discrimination. We emphasize that this in-distribution calibration is established on SIPaKMeD itself; whether it transfers to independent datasets is precisely the question taken up in the cross-dataset analysis, and the answer there is markedly less favourable.